\setlength{\parindent}{0pt}

\chapter{Com la química va canviar per sempre}

\epigraf{Cuanto peor, mejor para todos, y cuanto peor para todos, mejor, mejor para mí el suyo, beneficio político.}{Mariano Rajoy}

\lettrine[loversize=0.2, lines=2]{A}{ }la Facultat de Física de la Universitat dels Sabers Universals, on els estudiants només pensaven a arribar a entendre els misteris de l'univers, hi havia una assignatura que ningú volia fer: Química per a Físics. No era perquè la química fos una matèria misteriosa o fascinant, sinó perquè els futurs físics no veien cap sentit a estudiar-la. Els laboratoris, plens de tubs d'assaig, provetes i líquids de colors, els semblaven una pèrdua de temps comparat amb les equacions elegants i les teories de l'espai-temps.\\

Cada any, els nous estudiants de física entraven il·lusionats, somiant a descobrir les lleis fonamentals de l'univers, però ràpidament es trobaven atrapats en aquesta assignatura incòmoda. El temari no era res del que esperaven. No era una introducció lleugera a la química, ni tan sols un repàs dels principis bàsics. No, el que s'ensenyava era termodinàmica, una branca de la física que, a ulls dels alumnes de primer any, no tenia ni cap ni peus.\\

—Però què és això?! —es queixaven des de l'inici— Venim a estudiar física de debò, i en lloc d'entendre com funcionen les estrelles o els quarks, ens fan aprendre fórmules de calor, pressió i energia com si això ens importés.\\

Els estudiants passaven hores mirant gràfiques de cicles i intentant desxifrar equacions que no els explicaven gairebé res del món que volien comprendre. Fins i tot els professors de cursos superiors es posaven les mans al cap.\\

—Com pot ser que aquests alumnes arribin a la meva classe sense saber ni tan sols calcular correctament un camp elèctric? – es queixava el sergent Vall de Nuria, que ensenyava electromagnetisme.\\

—És clar —responia la professora Bagana, que donava mecànica clàssica— Si els teniu tot el primer any perdent el temps amb coses que no tenen cap context, és normal que es perdin!\\

El malestar entre els professors va anar creixent, fins que un dia, es van reunir en una gran sala, amb llums fluorescents que no paraven de fer pampallugues (ja ningú es molestava a arreglar-les, perquè, com deien a la facultat, "la física no arregla bombetes"). Era l'hora de prendre una decisió.\\

—Hem de canviar aquest temari —va dir en veu alta el professor Fliggs— La termodinàmica no té sentit per a aquests estudiants. Arriben a la classe i estan més perduts que mai.\\

Després de moltes hores de discussions, van arribar a una conclusió: el temari s'havia de canviar radicalment. Però, en comptes de simplificar-lo o posar-hi química bàsica que els alumnes poguessin relacionar amb altres assignatures, van prendre una decisió encara més estranya. Van decidir que en lloc de termodinàmica, els estudiants de primer any estudiarien... quàntica.\\

Quan la notícia va córrer per la facultat, els alumnes no s'ho podien creure. D'acord, la termodinàmica no era divertida, però la quàntica era una cosa que, fins i tot en les converses de passadís, semblava inabastable.\\

—Què volen, que aprenguem a fer màgia? —deia en Biel, un dels alumnes més escèptics— Si ni tan sols entenem com funciona una reacció química, com esperen que entenguem la probabilitat de trobar un electró en dos llocs alhora?\\

I així va començar un nou any acadèmic amb el nou temari. Els professors van entrar a classe amb els seus retoladors i els seus diagrames de funcions d'ona, mentre els estudiants intentaven comprendre com una partícula podia ser una ona, o com un gat dins d'una caixa podia estar viu i mort alhora.\\

—Això té fins i tot menys sentit que la termodinàmica! —es lamentava la Clàudia— És com si, en lloc d'ensenyar-nos a caminar, ens volguessin fer volar sense ales!\\

Al final, molts es van rendir. Els passadissos estaven plens de comentaris sarcàstics sobre com "tot és quàntic", i la química, aquella pobra assignatura, quedava cada vegada més oblidada i menyspreada.\\

—A qui li importa com reaccionen les molècules? —es deien entre ells— Nosaltres som físics, no químics!\\

Però el problema persistia: cada vegada més alumnes arribaven a segon curs sense tenir una base sòlida de física ni química. L'enigma del temari continuava sent un malson per a estudiants i professors. I tot i els canvis aparentment "modernitzadors", la sensació que la química no servia per a res anava creixent entre els físics.\\

I així, la Facultat de Física es va convertir en un lloc on ningú entenia per què s'ensenyava el que s'ensenyava. La termodinàmica, la quàntica i la química continuaven lluitant entre si, en una guerra sense fi per trobar el seu lloc en un món que, segons els estudiants, havia perdut completament el sentit.

\setlength{\parindent}{20pt}